% #714 图示：夹注宽度取整到夹注字宽的整数倍，为奇数时正文错位半个正文字
% 横排呈现（报告者用整页旋转实现竖排，jiazhu 内部始终在横排坐标系工作），
% 在正文下方标出 20pt 正文字网格，并画出夹注结束后正文的实际起点。
\documentclass{standalone}
\usepackage{xeCJK}
\usepackage{jiazhu}
\usepackage{xcolor}
\usepackage{tikz}
\setCJKmainfont{Noto Serif CJK SC}
\setmainfont{DejaVu Sans}
\parindent=0pt

\jiazhuset{beforeskip=0pt,afterskip=0pt,ratio=0.5}

% 测量「甲乙丙 + 夹注」的宽度，据此画出正文恢复处的位置刻度
\newlength\jzw
\newcommand\gridrow[2]{%
  \settowidth\jzw{\fontsize{20pt}{20pt}\selectfont 甲乙丙\jiazhu{#1}}%
  \begin{tikzpicture}[baseline=0pt]
    % 网格线画在内容下方，不被字形遮挡
    \foreach \i in {0,...,13}
      \draw[black!22,line width=0.4pt] (\i*20pt,-13pt) -- (\i*20pt,20pt);
    % 20pt 刻度标签
    \foreach \i in {0,2,4,6,8,10,12}
      \node[black!35,font=\tiny,anchor=north] at (\i*20pt,-13pt) {\i};
    % 内容
    \node[inner sep=0pt,anchor=south west] at (0,0)
      {\hbox{\fontsize{20pt}{20pt}\selectfont 甲乙丙\jiazhu{#1}丁戊}};
    % 夹注结束、正文恢复处的竖标记
    \draw[#2,line width=1.2pt] (\jzw,-13pt) -- (\jzw,20pt);
  \end{tikzpicture}}

\begin{document}
\begin{minipage}{500pt}
{\large\bfseries jiazhu \#714：夹注宽度只按夹注字宽取整，正文脱离字网格}

\vspace{3pt}
{\footnotesize 正文 20pt，\texttt{ratio=0.5} 故夹注字宽 10pt。灰线为 20pt 正文字网格（刻度以正文字为单位），
粗竖线标出夹注结束、正文恢复处。}

\vspace{10pt}
\gridrow{一二三四}{green!50!black}\\[1pt]
{\footnotesize\textcolor{black!60}{$n=4$：盒宽 $=2$ 个夹注字宽 $=1$ 个正文字宽}\quad\textcolor{green!45!black}{\textbf{正文恢复于刻度 4，落在网格上}}}

\vspace{9pt}
\gridrow{一二三四五}{red!75!black}\\[1pt]
{\footnotesize\textcolor{black!60}{$n=5$：盒宽 $=3$ 个夹注字宽 $=1.5$ 个正文字宽}\quad\textcolor{red!70!black}{\textbf{正文恢复于刻度 4.5，错位半个正文字}}}

\vspace{9pt}
\gridrow{一二三四五六}{red!75!black}\\[1pt]
{\footnotesize\textcolor{black!60}{$n=6$：盒宽 $=3$ 个夹注字宽}\quad\textcolor{red!70!black}{\textbf{同样错位半个正文字}}}

\vspace{9pt}
\gridrow{一二三四五六七八}{green!50!black}\\[1pt]
{\footnotesize\textcolor{black!60}{$n=8$：盒宽 $=4$ 个夹注字宽 $=2$ 个正文字宽}\quad\textcolor{green!45!black}{\textbf{正文恢复于刻度 5，落在网格上}}}

\vspace{12pt}
{\footnotesize 规律：盒宽 $=\lceil n/2\rceil$ 个夹注字宽。$\lceil n/2\rceil$ 为奇数（即 $n\bmod 4\in\{1,2\}$）时，
盒宽是正文字宽的半整数倍，其后所有正文都偏移半个字。}
\end{minipage}
\end{document}
